\documentclass[11pt]{article}
\usepackage[english]{babel}
\usepackage[a4paper,margin=0.9in]{geometry}
\usepackage[T1]{fontenc}
\usepackage[utf8]{inputenc}
\usepackage{lmodern}
\usepackage{amsmath,amssymb,mathtools}
\usepackage{microtype}
\usepackage{fancyhdr}
\usepackage{enumitem}
\setlength{\parindent}{1.5em}
\setlength{\parskip}{0.18em}
\emergencystretch=8em
\setlength{\headheight}{14pt}
\pagestyle{fancy}
\fancyhf{}
\lhead{Weber Vol. I \S\S 183--188}
\rhead{English translation}
\cfoot{\thepage}
\newcommand{\sect}[2]{\section*{\S\,#1. #2}}
\newcommand{\eps}{\varepsilon}
\newcommand{\Om}{\Omega}
\newcommand{\R}{\mathbb R}
\begin{document}
\begin{center}
{\large Volume I.}\par
{\Large Eighteenth Section.}\par\medskip
{\Large Roots of Metacyclic Equations.}\par\medskip
{\large \S\S\,183--188}
\end{center}

\sect{183}{Statement of the problem. Auxiliary theorem}

In the preceding section we found general criteria by which one can decide whether a proposed equation of prime degree is metacyclic. This by no means exhausts the subject. In the form in which Abel posed the problem, one must give a procedure by which all metacyclic equations can be found. Kronecker took up this problem and solved it completely, in the sense that the roots are found first, rather than the equations themselves.

For every prime number $n$, Kronecker gave an expression obtained from a given field $\Om$ by repeated extraction of roots. This expression has the following double property: every root of an irreducible metacyclic equation of degree $n$ in $\Om$ is contained in it, and conversely every such expression satisfies an irreducible equation of degree $n$ in $\Om$. The basis of the investigation is supplied by Lagrange's resolvents.

Let $n$ be a prime, let
\[
  x_0,x_1,x_2,\ldots,x_{n-1}
\]
be independent variables, numbered so that $x_n=x_0$, $x_{n+1}=x_1$, and so on. Let $\eps$ be an $n$-th root of unity. Put
\begin{equation}
  (\eps,x)=x_0+\eps x_1+\eps^2x_2+\cdots+\eps^{n-1}x_{n-1}
  =\sum_{h=0}^{n-1}\eps^h x_h .
\end{equation}
These are the Lagrange resolvents. Their behavior under cyclic permutations was studied in the fourteenth section. The point now is to determine the effect on them of the linear permutations considered in \S\,180.

We first recall an auxiliary theorem. The imaginary $n$-th roots of unity are the roots of $X=0$, where
\begin{equation}
  X=\xi^{n-1}+\xi^{n-2}+\cdots+\xi+1.
\end{equation}
This polynomial is irreducible over the rational field. If $\Om$ is a field over which $X$ remains irreducible and $\Om(\eps)$ is the field of rational functions of $\eps$ with coefficients in $\Om$, then each element of $\Om(\eps)$ can be written in one and only one way in the normal form
\begin{equation}
  c_0+c_1\eps+c_2\eps^2+\cdots+c_{n-2}\eps^{n-2},
\end{equation}
where the $c_i$ lie in $\Om$.

This remains true when $\Om$ is the field generated over the rational field by adjoining the independent variables $x_0,\ldots,x_{n-1}$. If $X$ became reducible there, one could specialize the variables $x_i$ to rational numbers without lowering the degree of the factors in $\eps$, and one would obtain a factorization over the rational field, contradicting irreducibility.

More generally, fix a permutation group $P$ on the $n$ indices. The rational functions in $x_0,
\ldots,x_{n-1}$ that are invariant under $P$ form a field, called the field belonging to $P$. Its elements are rationally expressible by the symmetric functions of the $x_i$ and by a function belonging to $P$. Thus the normal-form theorem applies to every such field.

Consequently, if a rational function $\Phi(x_0,x_1,\ldots,x_{n-1},\eps)$ in $\Om(\eps)$ remains unchanged when the permutations of a group $P$ are applied to the indices of the $x_i$, then the coefficients $c_0,c_1,\ldots,c_{n-2}$ in its normal form are themselves functions invariant under $P$.

If $\Phi(\eps)$ is an element of $\Om(\eps)$, its conjugates are obtained by the substitutions $(\eps,\eps^h)$, $h=1,2,\ldots,n-1$. Thus $\Om(\eps)$ is a normal field. If $\Phi(\eps)$ remains unchanged under all these substitutions, then it already lies in $\Om$. Indeed,
\[
  \Phi(\eps)=\frac1{n-1}\sum_{h=1}^{n-1}\Phi(\eps^h)
  =c_0+\frac{c_1+c_2+\cdots+c_{n-2}}{n-1},
\]
and the preceding assertion forces $c_1=c_2=\cdots=c_{n-2}=0$.

\sect{184}{Theorems on the resolvents}

We now consider the Lagrange resolvents $(\eps,x)$, where $\eps$ is an imaginary $n$-th root of unity and the $x_i$ are independent variables. It suffices to study the two generating substitutions
\begin{equation}
  s=(h,h+1),\qquad t=(h,gh),
\end{equation}
where the indices are taken modulo $n$ and $g$ is a primitive root modulo $n$.

Applying $s$ to
\begin{equation}
  (\eps,x)=x_0+\eps x_1+\eps^2x_2+\cdots+\eps^{n-1}x_{n-1}
\end{equation}
gives
\[
  x_1+\eps x_2+\eps^2x_3+\cdots+\eps^{n-1}x_0=\eps^{-1}(\eps,x).
\]
Applying $t$ replaces $x_h$ by $x_{gh}$ and gives
\begin{equation}
  x_0+\sum_{h=1}^{n-1}\eps^h x_{gh}
  =x_0+\sum_{h=1}^{n-1}\eps^{g^{-1}h}x_h=(\eps^{g^{-1}},x).
\end{equation}
Thus
\[
\begin{array}{ll}
(s): (\eps,x)\mapsto \eps^{-1}(\eps,x),&
(s^{-1}): (\eps,x)\mapsto \eps(\eps,x),\\[0.2em]
(t): (\eps,x)\mapsto (\eps^{g^{-1}},x),&
(t^{-1}): (\eps,x)\mapsto (\eps^g,x).
\end{array}
\]

Put
\begin{equation}
  (\eps,x)^\lambda=\Phi(\eps),\qquad
  (\eps^\lambda,x)(\eps,x)^{-\lambda}=F(\eps).
\end{equation}
When these functions are written in normal form, their coefficients are cyclic functions of $x_0,
\ldots,x_{n-1}$. With $\lambda=g$, define
\begin{equation}
\begin{aligned}
 f_0&=(\eps^g,x)(\eps,x)^{-g},\\
 f_1&=(\eps^{g^2},x)(\eps^g,x)^{-g},\\
 &\hspace{2em}\vdots\\
 f_{n-2}&=(\eps,x)(\eps^{g^{n-2}},x)^{-g}.
\end{aligned}
\end{equation}
Then $f_h=f(\eps^{g^h})$, with subscripts taken modulo $n-1$, and $t^{-1}$ cyclically permutes the $f_h$. Hence every cyclic function
\begin{equation}
  C(f_0,f_1,\ldots,f_{n-2})
\end{equation}
with rational coefficients is a metacyclic function of the variables $x_0,
\ldots,x_{n-1}$. Since the $f_h$ are also cyclically permuted when $\eps$ is replaced by $\eps^g$, such a cyclic function is independent of $\eps$.

Multiplying the defining equations for the $f_h$ after raising them to the powers $g^{n-2},g^{n-3},\ldots,1$ gives
\begin{equation}
  (\eps,x)^{g^{n-1}-1}=f_0^{g^{n-2}}f_1^{g^{n-3}}\cdots f_{n-2}.
\end{equation}
The primitive root $g$ may be chosen so that
\[
  {g^{n-1}-1\over n}\equiv1\pmod n.
\]
With
\begin{equation}
  \lambda={g^{n-1}-1\over n},\qquad
  F(\eps)=(\eps^\lambda,x)(\eps,x)^{-\lambda},
\end{equation}
we obtain
\begin{equation}
  (\eps,x)^n=F(\eps)^n f_0^{-g^{n-2}}f_1^{-g^{n-3}}\cdots f_{n-2}^{-1}.
\end{equation}
For the conjugate values of $\eps$, if
\begin{equation}
  F(\eps^{g^v})=F_v,
\end{equation}
then
\begin{equation}
  (\eps^{g^v},x)^n=F_v^n f_v^{-g^{n-2}}f_{v+1}^{-g^{n-3}}\cdots f_{v+n-2}^{-1}.
\end{equation}
Writing
\begin{equation}
  g^v=nq_v+r_v,\qquad 0<r_v<n,
\end{equation}
with $r_0=1$, and reducing the exponents to their least positive residues, one obtains functions $\Phi_v$ such that
\begin{equation}
  (\eps^{r_v},x)^n=\Phi_v f_v^{r_{n-2}}f_{v+1}^{r_{n-3}}\cdots f_{v+n-2}^{r_0}.
\end{equation}
The functions that occur here, denoted collectively by $\omega_v$, have four characteristic properties: they are fixed by the cyclic permutation $s=(h,h+1)$ of the indices of $x$; under $t^{-1}=(h,g^{-1}h)$ their indices are cyclically permuted; the substitution $(\eps,\eps^g)$ produces the same cyclic permutation; and cyclic functions of the $\omega_v$ are metacyclic functions of the $x_i$ and contain no $\eps$.

By Lagrange's theorem, any system $\Theta_0,\ldots,\Theta_{n-2}$ with these same four properties is rationally expressible in the field of metacyclic functions by one of the $\omega_v$, provided the $\omega_v$ are distinct. Indeed, if
\begin{equation}
  (u-
\omega_0)(u-
\omega_1)\cdots(u-\omega_{n-2})=\varphi(u),
\end{equation}
then $\varphi$ has metacyclic coefficients independent of $\eps$, and
\begin{equation}
  \chi(u)=\sum_{v=0}^{n-2}{\Theta_v\varphi(u)\over u-
\omega_v},\qquad
  {\chi(u)\over\varphi'(u)}=\Theta(u)
\end{equation}
gives $\Theta_v=\Theta(\omega_v)$.

\sect{185}{Roots of metacyclic equations}

We now substitute for $x_0,x_1,\ldots,x_{n-1}$ the roots $\xi_0,\xi_1,\ldots,\xi_{n-1}$ of an irreducible metacyclic equation of degree $n$ over $\Om$, in an order for which the metacyclic functions of the roots are rational in $\Om$. First assume that no resolvent $(\eps,\xi)$ vanishes and that no two of the functions $f_0,f_1,
\ldots,f_{n-2}$ become equal.

Then the $f$ become
\[
  k_0,k_1,\ldots,k_{n-2},
\]
the $F$ become
\[
  K_0,K_1,\ldots,K_{n-2},
\]
and the $k_i$ are the roots of a cyclic equation of degree $n-1$,
\begin{equation}
  \psi(u)=(u-k_0)(u-k_1)\cdots(u-k_{n-2})=0.
\end{equation}
Moreover,
\begin{equation}
  K_0=\Phi(k_0),\quad K_1=\Phi(k_1),\quad\ldots,\quad
  K_{n-2}=\Phi(k_{n-2}).
\end{equation}
Let
\begin{equation}
  \tau_v=\sqrt[n]{k_v}.
\end{equation}
Then
\begin{equation}
  (\eps^{r_v},\xi)=K_v\tau_v^{r_{n-2}}\tau_{v+1}^{r_{n-3}}
  \cdots\tau_{v+n-2}^{r_0},
\end{equation}
and, since $(1,\xi)=A$ is rational,
\begin{equation}
  n\xi_0=A+
  \sum_{v=0}^{n-2}K_v\tau_v^{r_{n-2}}\tau_{v+1}^{r_{n-3}}
  \cdots\tau_{v+n-2}^{r_0}.
\end{equation}
Although $n-1$ radicals occur, their values are not independent. With
\begin{equation}
  R_v=k_v^{r_{n-2}}k_{v+1}^{r_{n-3}}\cdots k_{v+n-2}^{r_0}
\end{equation}
one has
\begin{equation}
  K_v\sqrt[n]{R_v}=K_{v-1}^{g}\,k_{v-1}\,(\sqrt[n]{R_{v-1}})^{-g}.
\end{equation}
Thus all radicals can be expressed rationally through one of them and through the $k_i$.

Changing the determinations of the radicals gives
\begin{equation}
  n\xi=A+
  \sum_{v=0}^{n-2}E_vK_v\tau_v^{r_{n-2}}\tau_{v+1}^{r_{n-3}}
  \cdots\tau_{v+n-2}^{r_0},
\end{equation}
where the $E_v$ are $n$-th roots of unity determined by a single choice. Consequently the expression has only $n$ distinct values, namely the $n$ roots of the equation. In the natural cyclic order these roots are
\begin{equation}
  n\xi_h=A+
  \sum_{v=0}^{n-2}\eps^{-hr_v}K_v\tau_v^{r_{n-2}}\tau_{v+1}^{r_{n-3}}
  \cdots\tau_{v+n-2}^{r_0}.
\end{equation}

\sect{186}{Removal of the restricting hypotheses}

The hypotheses just imposed are not essential. Let $\eta_0,\eta_1,\ldots,\eta_{n-1}$ be the roots of any irreducible metacyclic equation of degree $n$. Introduce new unknowns
\begin{equation}
  \xi_0=\psi(\eta_0),\quad \xi_1=\psi(\eta_1),\quad\ldots,\quad
  \xi_{n-1}=\psi(\eta_{n-1}),
\end{equation}
where
\begin{equation}
  \psi(y)=a_0+a_1y+a_2y^2+\cdots+a_{n-1}y^{n-1}
\end{equation}
has coefficients in $\Om$. The coefficients $a_i$ can be chosen so that the $\xi_i$ are distinct and so that none of the required resolvents vanishes and no two of the resulting $k_i$ are equal. This follows because the transformation has Vandermonde determinant equal to the product of the differences of the $\eta_i$.

Writing conversely
\begin{equation}
  \eta_0=\chi(\xi_0),\quad \eta_1=\chi(\xi_1),\quad\ldots,\quad
  \eta_{n-1}=\chi(\xi_{n-1}),
\end{equation}
and setting
\begin{equation}
  { (\eps^v,y)\over (\eps^v,x)}=\Theta_v,
\end{equation}
one obtains the same four transformation properties as before. Hence
\begin{equation}
  Q_v=Q(k_v),\qquad
  (\eps^v,\eta)=Q_v(\eps^v,\xi),
\end{equation}
so that the same formula applies, with $K_v$ replaced by $Q_vK_v$.

The general theorem is therefore:
\begin{quote}
Every root $\xi$ of a metacyclic equation of prime degree $n$ can be represented as
\begin{equation}
  \xi=A+
  \sum_{v=0}^{n-2}K_v\tau_v^{r_{n-2}}\tau_{v+1}^{r_{n-3}}
  \cdots\tau_{v+n-2}^{r_0}.
\end{equation}
Here $A$ is rational, $k_0,k_1,\ldots,k_{n-2}$ are distinct non-zero roots of a cyclic equation of degree $n-1$, $\tau_v^n=k_v$, and $K_v$ is obtained from $k_v$ by one and the same rational function. The exponents $r_0,
\ldots,r_{n-2}$ are the least positive residues of $1,g,g^2,
\ldots,g^{n-2}$ modulo $n$, where $g$ is a primitive root modulo $n$.
\end{quote}
Conversely, every quantity of this form is a root of an equation of degree $n$ over $\Om$. It is irreducible except in the exceptional case where one root is rational and the remaining roots lie already in $\Om(\eps)$.

\sect{187}{Reality conditions}

Assume now that $\Om$ is a real field. By \S\,180 there are two kinds of metacyclic equations of prime degree $n$: those with one real root and $n-1$ imaginary roots, and those all of whose roots are real. Likewise cyclic equations of even degree have either all roots real or all roots imaginary.

If the cyclic equation whose roots are $k_0,k_1,\ldots,k_{n-2}$ has real roots, and if the radicals $\tau_v=\sqrt[n]{k_v}$ are chosen real, then the formula of \S\,186 shows that one obtains the case with one real root and the other roots imaginary. If the $k_v$ are imaginary, then $k_v$ and $k_{v+(n-1)/2}$ are conjugate. From the relations
\begin{equation}
  \sqrt[n]{R_{v+(n-1)/2}}=\Phi(k_v)
  \left(\sqrt[n]{R_v}\right)^{(n-1)/2},
\end{equation}
and
\begin{equation}
  \sqrt[n]{R_{v+(n-1)/2}}\sqrt[n]{R_v}=\Psi(k_v),
\end{equation}
one sees that the corresponding radicals are conjugate imaginary quantities; their products are real. Hence the summands occur in conjugate pairs and all the roots $\xi$ are real. Thus, for a real field $\Om$:
\begin{quote}
If the auxiliary cyclic equation for the $k_v$ has real roots, the corresponding metacyclic equation has one real root. If the auxiliary cyclic equation has imaginary roots, the corresponding metacyclic equation has all roots real.
\end{quote}

\sect{188}{Metacyclic equations of fifth degree}

The preceding formula gives a general representation of the roots of a metacyclic equation of prime degree. To make it explicit in degree $5$, one must describe the roots $k_0,k_1,k_2,k_3$ of a cyclic quartic equation, assuming they are distinct and non-zero.

Start with
\begin{equation}
  w=(k_0-k_2)(k_1-k_3).
\end{equation}
This never vanishes and is real over a real field. Under the cyclic permutation $(0,1,2,3)$ it changes sign, so put
\begin{equation}
  w=4\sqrt c.
\end{equation}
Then every function of the $k_i$ that changes sign under the cycle is a rational multiple of $\sqrt c$. Introduce rational quantities $a,b,c$ by
\begin{equation}
  (k_0-k_2)^2+(k_1-k_3)^2=8b,
  \qquad
  (k_0-k_2)^2-(k_1-k_3)^2=8a\sqrt c.
\end{equation}
They satisfy
\begin{equation}
  b^2=c(1+a^2),
\end{equation}
and
\begin{equation}
  k_0-k_2=2\sqrt{b+a\sqrt c},
  \qquad
  k_1-k_3=2\sqrt{b-a\sqrt c},
\end{equation}
with
\begin{equation}
  \sqrt c=\sqrt{b+a\sqrt c}\,\sqrt{b-a\sqrt c}.
\end{equation}
If
\begin{equation}
  k_0+k_1+k_2+k_3=4C,
  \qquad
  k_0-k_1+k_2-k_3=4B\sqrt c,
\end{equation}
then
\begin{equation}
\begin{aligned}
 k_0&=C+B\sqrt c+\sqrt{b+a\sqrt c},\\
 k_1&=C-B\sqrt c+\sqrt{b-a\sqrt c},\\
 k_2&=C+B\sqrt c-\sqrt{b+a\sqrt c},\\
 k_3&=C-B\sqrt c-\sqrt{b-a\sqrt c}.
\end{aligned}
\end{equation}
These are the roots of a cyclic quartic equation. With $k-C=x$, that equation is
\begin{equation}
  x^4-2(B^2c+b)x^2-4Bac\,x+B^4c^2-2B^2bc+c=0.
\end{equation}

A special case must first be separated off: if $b=0$, then $a=i$, and the formula gives a cyclic quartic only when $i$ belongs to the base field; it therefore never occurs over a real field.

To remove the relation among $a,b,c$, assume $b\ne0$ and set
\begin{equation}
  b=h(1+a^2),\qquad c=h^2(1+a^2).
\end{equation}
Replacing $Bh$ by $B$, one obtains
\begin{equation}
\begin{aligned}
 k_0&=C+B\sqrt{1+a^2}+\sqrt{h(1+a^2+a\sqrt{1+a^2})},\\
 k_1&=C-B\sqrt{1+a^2}+\sqrt{h(1+a^2-a\sqrt{1+a^2})},\\
 k_2&=C+B\sqrt{1+a^2}-\sqrt{h(1+a^2+a\sqrt{1+a^2})},\\
 k_3&=C-B\sqrt{1+a^2}-\sqrt{h(1+a^2-a\sqrt{1+a^2})}.
\end{aligned}
\end{equation}
Abel's equivalent expression for $k_0$ is\n\[\n  k_0=C+B\sqrt{1+e^2}+\sqrt{h(1+e^2+\sqrt{1+e^2})},\n\]\nand is obtained from the preceding formula by a change of notation.

For degree $5$ take $g=2$, so that
\[
  r_0=1,\qquad r_1=2,
  \qquad r_2=4,
  \qquad r_3=3.
\]
Writing
\[
  \tau_i=\sqrt[5]{k_i}
\]
for $i=0,1,2,3$, the root is
\begin{equation}
\begin{aligned}
 \xi={}&A+K_0\tau_0^3\tau_1^4\tau_2^2\tau_3
       +K_1\tau_1^3\tau_2^4\tau_3^2\tau_0\\
      &+K_2\tau_2^3\tau_3^4\tau_0^2\tau_1
       +K_3\tau_3^3\tau_0^4\tau_1^2\tau_2.
\end{aligned}
\end{equation}
The coefficients $K_0,K_1,K_2,K_3$ are cyclically permuted rational functions of the $k_i$. Write
\begin{equation}
  r=\sqrt{1+a^2},\qquad
  \rho=\sqrt{h(1+a^2+a\sqrt{1+a^2})},
\end{equation}
\begin{equation}
  \rho'=\sqrt{h(1+a^2-a\sqrt{1+a^2})},
  \qquad \rho\rho'=hr.
\end{equation}
Then, for rational $A_1,A_2,A_3,A_4$,
\begin{equation}
\begin{aligned}
 K_0&=A_1+A_2r+A_3\rho+A_4r\rho,\\
 K_1&=A_1-A_2r+A_3\rho'-A_4r\rho',\\
 K_2&=A_1+A_2r-A_3\rho-A_4r\rho,\\
 K_3&=A_1-A_2r-A_3\rho'+A_4r\rho'.
\end{aligned}
\end{equation}
These four quantities are themselves the roots of a cyclic quartic equation. They are only apparently more general than the $k_i$, since
\[
  K_0=A_1+A_2r+\sqrt{\rho^2(A_3+A_4r)^2}
\]
has the same form as the expression for $k_0$.

\section*{Corrections}

Page 182, in the formula on line 12, read $x_m$ instead of $x_n$.

Page 347, in the formula on line 24, read $(2x^2+1)^2$ instead of $(2x^2-1)^2$.

\end{document}
